\section{Consecutive interlacing through the two-cycle range}
\label{sec:two-cycle-insertion}

\textbf{Two-cycle interlacing theorem.}
Let $k\ge3$. For every $k\le n\le3k-2$, put $P_n(t)=E_{n,k}(t)/t$.
Then $P_n$ and $P_{n+1}$ have simple negative zeros, and the zeros of
$P_{n+1}$ strictly interlace those of $P_n$ with the new polynomial
supplying both outer zeros. More precisely, at every zero $x$ of $P_n$,
\begin{equation}
 \left|L(x)\right|<\left|D(x)\right|,\qquad
 D(t)=t(1-t)P_n'(t),\qquad
 L(t)=\binom{n}{k-1}A_{k-1}(t)E_{n+1-k,k}(t).
 \label{eq:two-cycle-lag}
\end{equation}
Thus the conclusion covers every consecutive pair for which the new row
admits at most two cycles.

To obtain the insertion recurrence, either insert the largest label
into an edge of an existing cycle or place it in a new $k$-cycle.
Inserting it after the initial label of an ascending edge preserves
the excedance count. Inserting it in any other edge adds one
excedance. With the convention $E_{0,k}=1$ and $E_{m,k}=0$ for
$1\le m<k$, this gives
\begin{equation}
 P_{n+1}(t)=\bigl(1+(n-1)t\bigr)P_n(t)+D(t)+L(t).
 \label{eq:two-cycle-recurrence}
\end{equation}
The two polynomials $P_{n+1}$ and $P_{n+1}-2L$ therefore take the values
$D+L$ and $D-L$, respectively, at an old root. For $k\le n\le2k-2$
only single cycles occur in both rows and $L=0$. Eulerian
interlacing follows from
$A_{m+1}=(1+mt)A_m+t(1-t)A_m'$ by induction from $A_1=1$,
evaluating the recurrence at the simple negative zeros of $A_m$.
The first nonzero insertion term, at $n=2k-1$, is covered by the
\href{https://github.com/ProgramComputer/open-math-research/pull/27}{reviewed first-transition argument}.
It compares $A_{2k-2}$ with both polynomials
$A_{2k-1}\pm\binom{2k-1}{k-1}tA_{k-1}^2$.
We treat $2k\le n\le3k-2$ below.

\subsection{The first ranks and their derivative bound}

Write $M=n-1$, $E=E_{n,k}$, and
$\alpha_p=[t^p]E/p^M$ for $p\ge1$. We take $n\ge100$ in this
subsection. The bounds needed for the first ranks are
\begin{equation}
 0\le\alpha_p\le3\quad(p\ge1),\qquad
 \alpha_p\ge99/100\quad
 \left(p\le \frac{451}{1000}\sqrt M\right).
 \label{eq:two-cycle-coefficients}
\end{equation}
A permutation with $p-1$ descents gives a word on $p$ colors by assigning
each label to its increasing run. The word determines the permutation,
so $[t^p](tA_m)\le p^m$. Conversely, sorting each color class and
concatenating the classes in color order creates all $p-1$ descents
if every class meets both label halves. The union bound shows that
the single-cycle term $tA_M$ alone gives
$\alpha_p\ge 1-2p\exp(-\lfloor M/2\rfloor/p)$. For the stated
range of $p$, this error is below
$\sqrt M\exp(-9\sqrt M/10)<1/100$.

Each unordered pair of cycles is counted twice in the sum
\begin{equation}
 E(t)=tA_{n-1}(t)+\frac12\sum_{j=k}^{n-k}\binom nj
                (tA_{j-1}(t))(tA_{n-j-1}(t)).
 \label{eq:two-cycle-exact-polynomial}
\end{equation}
Consequently the coefficient bound from words, followed by symmetry
under $a\mapsto p-a$, gives
\[
 \alpha_p\le
 1+\frac12\sum_{a=1}^{p-1}\frac{p}{a(p-a)}
        \Pr\{\operatorname{Bin}(n,a/p)\in[k,n-k]\}
 =
 1+\sum_{a=1}^{p-1}\frac1a
        \Pr\{\operatorname{Bin}(n,a/p)\in[k,n-k]\}.
\]
Since $k\ge n/3$, Hoeffding's inequality bounds the probability by
$\exp(-8n/225)$ whenever $a/p\notin(1/5,4/5)$. For $p\le n-1$ the sum of
these terms is at most $(1+\log n)\exp(-8n/225)<1/6$.
The expression decreases for $n\ge100$ and satisfies the bound at
$n=100$ using $\log100<14/3$ and $\exp(32/9)>34$.
For the remaining indices the sum of $1/a$ is at most $11/6$.
When $p\le4$, this follows by summing through $a=3$.
When $5\le p\le11$, all the indices lie in $[2,8]$.
For $p\ge12$, integrate $1/a$ from the first index to the last,
allowing the first term separately, to obtain
$\log4+5/p\le\log4+5/12<11/6$. Coefficients beyond $n-1$ vanish.
This proves \eqref{eq:two-cycle-coefficients}.

Put $H=\lfloor(7/20)\sqrt M\rfloor$.
For a rank $1\le h\le H$ define
\[
 \rho=\frac{\alpha_h}{\alpha_{h+1}},\qquad
 u=\rho\left(\frac h{h+1}\right)^M,\qquad
 G(z)=\frac{E(uz)}{\alpha_h h^M(uz)^h},\qquad
 b_p=\left[\frac ph\left(\frac h{h+1}\right)^{p-h}\right]^M.
\]
We have $h+1\le(451/1000)\sqrt M$ and
$\rho,\rho^{-1}\le R=100/33$. The terms of $G$ with indices
$h,h+1$ are exactly $1,z$. The successive ratios of $b_p$ away from
this pair are at most $(1-1/h^2)^M$ on the left and
$(1-1/(h+1)^2)^M$ on the right. Both are bounded by
$q_0=\exp(-1000000/203401)<1/135$, and the left side is empty
when $h=1$.

Whenever $q_0\max(|\rho z|,|\rho z|^{-1})\le\eta$, the absolute values
of the terms of $G$ at indices $h-\ell$ and $h+1+\ell$ are at most
$R\eta^\ell$ and $R|z|\eta^\ell$, respectively, for $\ell\ge1$.
The second estimate uses the coefficient ratio to $\alpha_{h+1}$,
as the right central term equals $z$.
On $|z+1|=1/4$ we may take $\eta=3/100$, since
$3/4\le |z|\le5/4$. Here
\[
 |G(z)-(1+z)|\le R(1+|z|)\frac{\eta}{1-\eta}<1/4.
\]
On $|z|=1/2$ and $|z|=3$ we may take $\eta=7/100$, and the same
bound is below $1/2$ and $2$, respectively. Rouch\'e's theorem
applied to $z^{h-1}G(z)$ gives $h-1$ roots inside $|z|=1/2$, $h$
inside $|z|=3$, and one in $|z+1|<1/4$. That disk lies in the
annulus between the two circles. Its sole root is real by
conjugation. Thus $P_n$ has a simple negative root $x=uz$ of absolute
rank $h$ with $3/4<|x|/u<5/4$.

At this root, put $X=|x|$. The normalized derivative
$xE'(x)/(\alpha_h h^M x^h)$ equals $zG'(z)-G(z)$ because
$G(z)=0$. Its central terms contribute $-1$ together. The remaining
weights are $\ell+1$ and $\ell$ on the left and right. With
$\eta=3/100$ this gives
\begin{equation}
 \begin{split}
 \bigl|zG'(z)-G(z)+1\bigr|
 &\le R\left((1+5/4)\frac{\eta}{(1-\eta)^2}
                    +\frac{\eta}{1-\eta}\right)<\frac13,\\
 |xE'(x)|&>\frac23\alpha_h h^M X^h .
 \end{split}
 \label{eq:two-cycle-front-derivative}
\end{equation}

The residual row has size $n+1-k$ between $k$ and $2k-1$.
It consists of single cycles, so
$L(t)=\binom{n}{k-1}tA_{k-1}(t)A_{n-k}(t)$. Applying the same
coefficient bound to both Eulerian factors yields
\[
 [t^r](tL)
 \le \frac{n}{n-k+1}\,r^M
     \sum_{a=1}^{r-1}
         \Pr\{\operatorname{Bin}(M,a/r)=k-1\}.
\]
Here $b=(k-1)/M$ belongs to $[1/3,1/2]$, and Stirling's bounds give
\[
 \Pr\{\operatorname{Bin}(M,b)=k-1\}
 \le \frac{\exp(1/(12M))}{\sqrt{2\pi M b(1-b)}}
 < \frac{7}{8\sqrt M}.
\]
Changing the binomial parameter from $b$ to $y$ multiplies this
point probability by $\exp(-M I(b,y))$, where
$I(b,y)=b\log(b/y)+(1-b)\log((1-b)/(1-y))$.
As a function of $b$, $I$ and its first derivative vanish at $b=y$,
and its second derivative is at least four. Thus
$I(b,y)\ge2(b-y)^2$. The Gaussian sum on the lattice $y=a/r$ is
at most $1+\sqrt{\pi/2}\,r/\sqrt M$. Indeed, each superlevel
interval contains at most its length plus one lattice points in
the $a$ coordinate, and integrating the level from zero to one
gives the bound. Since $n/(n-k+1)<2$,
\begin{equation}
 [t^r](tL)\le \frac{2(7/8)}{\sqrt M}\,
       r^M\left(1+\sqrt{\pi/2}\,\frac r{\sqrt M}\right).
 \label{eq:two-cycle-front-lag}
\end{equation}
For the sum at $x$, the weights
$r^M X^r/(\alpha_h h^M X^h)$ are bounded by
$(100/99)\eta^\ell$ at $r=h-\ell$ and
$(100/99)(5/4)\eta^\ell$ at $r=h+1+\ell$, for $\ell\ge0$.
For the right-hand weights use $1/\alpha_{h+1}$, which absorbs
the factor $\rho$ from the first step. The weight in parentheses in
\eqref{eq:two-cycle-front-lag} is at most $157/100$ through
$r=h+1$ and grows by less than $13/100$ per step thereafter. The
normalized weighted sum is therefore below
\[
 \frac{100}{99}\left(\frac94\frac{157}{100(1-\eta)}
          +\frac54\frac{13\eta}{100(1-\eta)^2}\right)
 <\frac{37}{10}.
\]
Since $E'(x)=xP_n'(x)$, equations
\eqref{eq:two-cycle-front-derivative} and
\eqref{eq:two-cycle-front-lag} give
\begin{equation}
 \frac{|L(x)|}{|D(x)|}
 =\frac{|xL(x)|}{(1+X)|xE'(x)|}
 <\frac{777}{80\sqrt M}<1.
 \label{eq:two-cycle-front-ratio}
\end{equation}

\subsection{Middle-root signs and a fixed rational comparison}

For $t=-e^{-s}$, $s\ge0$, let
\[
 r=\sqrt{s^2+\pi^2},\qquad
 \theta=\arg(s+i\pi),\qquad S=\sin\theta,\qquad
 z_p=(\cos\theta+ipS)^{-1},\qquad
 B_m=\sum_{p\ \mathrm{odd}}z_p^m .
\]
Here $p$ runs over both positive and negative odd integers.
The Eulerian power-sum identity gives the normalized $m$-cycle weight
\begin{equation}
 \frac{r^m tA_{m-1}(t)}{m!(1-t)^m}=\frac{B_m}{m}\qquad(m\ge2).
 \label{eq:two-cycle-pole}
\end{equation}
Indeed $tA_{m-1}(t)/((m-1)!(1-t)^m)$ is the coefficient of $w^{m-1}$
in $1/(1+e^{w-s})-1$. The identity at $t=-1$ follows by continuity.
The poles $s+ip\pi$ all have residue $-1$.
For fixed $s$, the circles of radii $\sqrt{s^2+4b^2\pi^2}$, with integers
$b\to\infty$, stay a positive distance from the poles. On these
circles the reciprocal denominator is bounded: if the real part
of $w-s$ is bounded away from zero the reverse triangle inequality
applies, and otherwise its imaginary part stays away from the odd
multiples of $\pi$. The coefficient integral on the circle tends to
zero. Summing the residues in the coefficient formula proves
\eqref{eq:two-cycle-pole}, since the sum converges absolutely.

Define
\[
 F_n=\frac{r^n E_{n,k}(t)}{n!(1-t)^n},\qquad
 f=\frac n2F_n
   =\frac{B_n}{2}+\sum_{j=k}^{n-k}w_j\frac{B_jB_{n-j}}{2},
 \qquad w_j=\frac{n}{2j(n-j)}.
\]
The identity follows by applying \eqref{eq:two-cycle-pole} to
\eqref{eq:two-cycle-exact-polynomial}. Set $W=\sum_{j=k}^{n-k}w_j$.
Then $W=\sum_{j=k}^{n-k}1/j\le\log 2<7/10$, by integration from
$k-1$ and the bound $n-k\le2k-2$.

Assume $n\ge N=192$ and $nS^2\ge c=6/5$.
Let $\beta=c/N=1/160$ and $K(p)=1+\beta(p^2-1)$. If $m/n\ge a$,
the monotonicity of $(1+u/n)^n$ in $n$ for $u\ge0$ gives
$|z_p|^m\le K(p)^{-H_a}$, where $H_a=Na/2$.
All the parameters $a$ used below are at least $1/3$ and are integer
multiples of $1/48$, so $H_a$ is an integer. With
\[
 I_a=\frac{K(9)^{1-H_a}}{4\beta(H_a-1)},\qquad
 T(a)=\sum_{p=3,5,7,9} K(p)^{-H_a}+\frac{I_a}{9},\qquad
 U(a)=\sum_{p=3,5,7,9}pK(p)^{-H_a}+I_a ,
\]
we obtain, for $e_m=B_m/2-\cos(m\theta)$,
\begin{equation}
 |e_m|\le T(a),\qquad |e_m'|\le mU(a).
 \label{eq:two-cycle-single-tails}
\end{equation}
The term with $|p|=1$ gives the displayed cosine. For the derivative,
$|z_p'|\le |p|\,|z_p|^2$ and $|z_p|\le1$. The sum on positive odds
above nine is at most half the integral starting at nine of the
decreasing majorant. The weighted integrand $pK(p)^{-H_a}$ decreases
there because $(2H_a-1)\beta p^2>1-\beta$. Its half integral equals
$I_a$, and the unweighted integrand is at most one ninth of the weighted
one on that interval. These summable bounds also justify
differentiating the series in this region.

Suppose $u\le\min(j,n-j)/n\le v$ for one adjacent mesh interval
$[u,v]$. Put
\begin{align*}
 A_{uv}&=1+2\bigl(T(u)+T(1-v)+T(u)T(1-v)\bigr),\\
 C_{uv}&=1-2u+
 2v\bigl(T(1-v)+U(u)(1+T(1-v))\bigr)\\
 &\hspace{2.65em}
 +2(1-u)\bigl(T(u)+U(1-v)(1+T(u))\bigr).
\end{align*}
For the two principal cosines, the identity
$2\cos(j\theta)\cos((n-j)\theta)=\cos(n\theta)+\cos((2j-n)\theta)$
gives a residual bounded by one in magnitude and $n(1-2u)$ in
derivative magnitude. Applying \eqref{eq:two-cycle-single-tails} and
the product rule to the terms involving at least one $e_m$ bounds
the full residual by $A_{uv}$ and its derivative by $nC_{uv}$.

For the threshold interval with index $i\in\{16,\ldots,23\}$ choose
$q=i/48\le k/n\le q_+=(i+1)/48$. Let $A_i,C_i$ be the respective
maxima of $A_{uv},C_{uv}$ over the adjacent mesh intervals in
$[q,1/2]$. With $T_*=T(1)$ and $U_*=U(1)$ this proves
\begin{equation}
 \begin{split}
 |f-(1+W)\cos(n\theta)|&\le T_*+WA_i,\\
 |f'/n+(1+W)\sin(n\theta)|&\le U_*+WC_i.
 \end{split}
 \label{eq:two-cycle-mixed-tails}
\end{equation}

To compare the insertion terms at a root $t=x$, multiply $D(x)$
and $L(x)$ by $x r^{n+1}/(n!(1-x)^{n+1})$.
Their transformed values are, respectively,
$S F_n'(\theta)$ and $B_kB_{n+1-k}/(n+1-k)$.
For the first equality use $t\,d\theta/dt=\pi/r^2$.
Terms from differentiating the normalization vanish at a zero of the
old row. For the second, apply
\eqref{eq:two-cycle-pole} to each of the two factors of $tL$.
Dividing further by $2S$ gives derivative magnitude $|f'|/n$
and lag magnitude at most
\[
 J_i=\frac{2(1+T(q))(1+T(1-q_+))}{15(1-q_+)},
\]
since $(n+1-k)S\ge(1-q_+)\sqrt{nc}>15(1-q_+)$.
The weight $W$ satisfies $0\le W\le W_i$ with
\[
 W_i=\min\left(\frac7{10},\frac{1-2q+1/N}{q}\right).
\]
The second bound uses at most $n-2k+1$ terms each at most $1/k$.

Define
\[
 R_i(w)=\frac{(T_*+wA_i)^2+(U_*+wC_i+J_i)^2}{(1+w)^2}.
\]
The table gives upper integer bounds for $10^6R_i(w)$ at the two
endpoints, computed by rational arithmetic from the formulas above.
The last two columns are ceilings.
\[
\begin{array}{c|c|r|r}
 i & W_i & \lceil 10^6R_i(0)\rceil & \lceil 10^6R_i(W_i)\rceil\\ \hline
 16 & 7/10  & 85473  & 988936\\
 17 & 7/10  & 88476  & 921937\\
 18 & 49/72 & 92155  & 844623\\
 19 & 41/76 & 96575  & 649634\\
 20 & 33/80 & 101824 & 488795\\
 21 & 25/84 & 108020 & 358384\\
 22 & 17/88 & 115311 & 256385\\
 23 & 9/92  & 123892 & 182547
\end{array}
\]
Since $A_i\ge1$, the polynomial
\[
 (T_*+wA_i)^2+(U_*+wC_i+J_i)^2-(1+w)^2
\]
is convex in $w$. Its endpoint values are strictly negative, so
\begin{equation}
 (T_*+WA_i)^2+(U_*+WC_i+J_i)^2<(1+W)^2.
 \label{eq:two-cycle-quadratic}
\end{equation}
At a checkpoint $\theta=l\pi/n$ in this region,
\eqref{eq:two-cycle-mixed-tails} and \eqref{eq:two-cycle-quadratic}
give $\operatorname{sgn} f=(-1)^l$.
At a zero of $f$, the first bound in \eqref{eq:two-cycle-mixed-tails}
bounds the cosine in the same inequality. Hence
$(1+W)|\sin(n\theta)|>U_*+WC_i+J_i$. The second bound in
\eqref{eq:two-cycle-mixed-tails} gives $|f'|/n>J_i$.
This proves \eqref{eq:two-cycle-lag} at all old roots in the
middle region on the side $-1\le t<0$.

\subsection{Joining the regions and counting signs}

Take $H=\lfloor(7/20)\sqrt{n-1}\rfloor$ as above and suppose $n\ge192$.
At the first checkpoint $\theta_{H+1}=(H+1)\pi/n$ we have
$n\sin^2\theta_{H+1}>6/5$. Indeed
$b=(7\pi/20)\sqrt{n-1}/n<\theta_{H+1}$ and
$b^2<121/(100n)$. Using $\sin b\ge b-b^3/6$ and
$333/106<\pi<22/7$ gives
\[
 n\sin^2\theta_{H+1}>
 \left(\frac7{20}\frac{333}{106}\right)^2
 \frac{191}{192}\left(1-\frac{121}{115200}\right)^2>\frac65.
\]
All the first $H$ roots precede this checkpoint in increasing modulus.
Their moduli are bounded by $(5/4)R(H/(H+1))^{n-1}$, and the
checkpoint corresponds to $s=\pi\cot\theta_{H+1}<n/(H+1)$.
With $M=n-1$, the positive Taylor series of
$-\log(1-1/(H+1))$ gives
\[
 M\log(1+1/H)-\frac{M+1}{H+1}
 \ge\frac{M}{2(H+1)^2}-\frac1{H+1}
 \ge\frac{1}{2(451/1000)^2}-\frac14
 >\log((5/4)R).
\]
This proves the strict separation.

Let $t_l=-\exp(-\pi\cot(l\pi/n))$ for
$l=H+1,\ldots,n-H-1$. The lower-half checkpoints through
$l=\lfloor n/2\rfloor$ satisfy the middle condition,
so $E_{n,k}(t_l)$ has sign $(-1)^l$ there. The upper-half
checkpoints follow by reciprocity. In fact
$t_{n-l}=1/t_l$ and $E_{n,k}(1/t)=t^{-n}E_{n,k}(t)$,
so the signs are again $(-1)^l$.
The $n-2H-2$ gaps between consecutive checkpoints each contain a
negative old root. The first $H$ roots and their $H$ reciprocal
partners lie strictly outside these gaps. This exhausts the degree
$n-2$ of $P_n$, so all old roots are simple.
For odd $n$ the center belongs to the central gap, and for even $n$
it is a checkpoint. Every old root on the lower side is therefore
covered by \eqref{eq:two-cycle-front-ratio} or
\eqref{eq:two-cycle-quadratic}. At an old root $x$, the reciprocal
identity for $P_n$ implies $D(1/x)=x^{1-n}D(x)$.
The cycle factors of $L$ give $L(1/x)=x^{1-n}L(x)$.
This transfers \eqref{eq:two-cycle-lag} to the upper side as well.

Since $P_n(0)=1$, the sign of $D$ at the old root of absolute rank
$h$ is $(-1)^h$. The two next-row polynomials $P_{n+1}$ and
$P_{n+1}-2L$ have constant and leading coefficients one, with
degree $n-1$. At every old root their signs equal the sign of $D$
by \eqref{eq:two-cycle-lag}. Both therefore have a root in each
successive gap and in each of the two outer intervals on the
negative axis. This exhausts their degree and proves strict
interlacing for $n\ge192$.

\subsection{Exact replay for the remaining pairs}

The remaining range has $3\le k\le95$ and
$2k\le n\le\min(191,3k-2)$. It contains $3\,007$ consecutive pairs
with $379\,502$ old roots, counting reciprocal partners and any
central roots. The finite verifier constructs the original coefficients
by distinguishing the cycle containing label one:
\[
 P_n(t)=A_{n-1}(t)+t\sum_{j=k}^{n-k}
              \binom{n-1}{j-1}A_{j-1}(t)A_{n-j-1}(t).
\]
It checks \eqref{eq:two-cycle-recurrence} separately as an integer
polynomial identity and checks both next-row polynomials against the
same old polynomial. Eulerian coefficients are reconstructed by their
recurrence and cross-checked by finite differences. A separate
Eulerian sign computation and the transition at $n=2k-1$ initialize the
brackets for each threshold.

Below modulus one, the verifier locates disjoint dyadic brackets
with denominator $2^{256}$. The signs of the old polynomial at the
bracket endpoints of rank $h$ are $(-1)^{h-1}$ and $(-1)^h$ in
order of increasing modulus. The brackets, their reciprocal partners,
and the center when the old degree is odd exhaust the old degree.
Each next-row polynomial is tested for sign $(-1)^h$ at both endpoints of that bracket. The sign
calculation uses scaled integer Horner evaluation. At each step the
rounding error is less than one scaled unit, and multiplication
by the modulus contracts the accumulated error. Hence an accumulator
with magnitude exceeding the degree in those units certifies the sign.
The endpoint signs at zero and $-1$ are evaluated exactly.

For each next-row polynomial, these signs force a lower root in
every gap between old brackets and one before the first bracket.
If the old degree is odd, the sign checked at its central root
forces an additional new root between the last lower bracket and
$-1$. If the new degree is odd, reciprocity gives its central root.
Reciprocal partners and these roots exhaust the new degree.
Thus each next-row polynomial has the same fixed sign throughout
each old bracket, and the central check and reciprocity give the
corresponding signs at the other old roots. This proves
\eqref{eq:two-cycle-lag} and strict interlacing for all finite pairs.

The \href{https://github.com/ProgramComputer/open-math-research/tree/research/topics/long-cycle-excedances}{topic archive}
contains the rational bound checker, the finite verifier and its replay
materials. Run the following from that topic directory with Python~3
and \texttt{python-flint} installed. FLINT is used for integer
polynomial arithmetic, and root signs use integer Horner evaluation.
\begin{verbatim}
python3 -I -O verify_two_cycle_bounds.py
python3 -I -O verify_two_cycle_band.py
\end{verbatim}
The second command reconstructs the brackets and needs no root archive.
Its optional \texttt{--records} argument writes a gzip JSONL index with
the certified pair ranges and hashes of the coefficients and endpoints.
The first command checks all sixteen rational endpoint comparisons and
the checkpoint join. No Lean or Rocq verification is asserted here.
